% =====================================================================
%  note.tex —— 译者说明（体例与使用方式）
% =====================================================================
\chapter*{译者说明}
\addcontentsline{toc}{chapter}{译者说明}
\markboth{译者说明}{}

本译注本是立川裕二（Yuji Tachikawa）为东京大学物理系研究生开设的“数理物理学”课程讲义
\textit{Algebraic topology for physicists}（2025 年 7 月稿，共 206 页）的中文译本，分两部分：
第一部分为第一学期材料（§1--§8），第二部分为第二学期补充材料（§9--§13）。

\para{页码对应。}底本有连续印刷页码，原书页与 PDF 页一一对应。译文中每个原书页起始处都有
注释 \verb|% [原文 p.N]| 作定位锚点；正文中引用的页码、公式号、定义/定理号一律照原文。

\para{数学与术语。}译文以“数学正确性优先于一切”为第一原则：原文的每一个定义、定理、
公式均回原页面图像逐符号核对，抽取文字层只作初稿。术语以书末《术语对照表》为准；
表中无稳定中文译名者保留英文原文。所有专有名词、人名不音译。

\para{译注体例。}本译本另加两类教学层，均可由导言区开关整体关闭：

\begin{itemize}
  \item 每节末的“本节重点”框：概括该节做了什么、结论是什么、为什么重要（40--90 字）；
  \item 每章末的“本章回顾与评注”框：只概括本章主线与各节推进方式，不评价、不补充原文（150--250 字）。
\end{itemize}

正文中另有行内批注 \verb|% [Tutor Note: ...]|，不占版面，用于三类情形：原文笔误或记号冲突、
跨节依赖、学生常见的易错点。原文笔误一律登记在交付说明（README）的“原文笔误登记表”中。
译文保留原文的灰底陈述框、编号公式与图号，不重新编号。

\para{参考文献与致谢。}参考文献按学术惯例不翻译，照原文转录英文原貌（个别条目中的日文作者名
保留原字）；作者信息、致谢按原文译出。原著中的引用键（如 [三24]）原样保留。
